\section*{ID10245: Definition of λ}
\paragraph{Metadata.}
PiID: 7852902345292 | RTEID: RTE:P39783.6102:T2.7193 | UUID: 25c632d9-549a-5088-88ca-7d18bb22e4db

\paragraph{Canonical Statement.}
Trawin Topology introduces constrained dipole dynamics and toroidal flux topology as fundamental organizing principles, providing a mathematical bridge between discrete quantum phenomena and continuous classical fields. We present rigorous proofs demonstrating wavelength quantization \textbackslash{}[ \textbackslash{}lambda = \textbackslash{}pi R\_\textbackslash{}oplus \textbackslash{}approx 2.0015\textbackslash{}times 10\textasciicircum{}\{4\}\textbackslash{},\textbackslash{}mathrm\{km\}, \textbackslash{}] tidal deformation parameters \textbackslash{}[ \textbackslash{}varepsilon \textbackslash{}approx 7.27\textbackslash{}times 10\textasciicircum{}\{-8\}, \textbackslash{}] and beat frequency mechanisms \textbackslash{}[ T\_\{\textbackslash{}text\{semidiurnal\}\} \textbackslash{}approx 12.42\textasciitilde{}\textbackslash{}text\{hours\}, \textbackslash{}] that validate the framework's predictions. This work positions Trawin Topology as a leading candidate for the geometric unification of fundamental forces.

\paragraph{Canonical Equation.}
\[
\lambda = \pi R_\oplus \approx 2.0015\times 10^{4}\,\mathrm{km},
\]

\paragraph{Dictionary.}
Trawin Topology introduces constrained dipole dynamics and toroidal flux topology as fundamental organizing principles, providing a mathematical bridge between discrete quantum phenomena and continuous classical fields.

\paragraph{Encyclopedia.}
Definition of λ is a scaling / order-of-magnitude relation, written λ = π R\_oplus ≈ 2.0015× 10\textasciicircum{}4 km. It expresses a scaling or proportionality, capturing how the quantity grows with its parameters. It is built from π, R, defining λ. Within the Trawin Zero Point registry it serves as one element of the framework's mathematical narrative.
